% TOP_PROBLEM: {"id": 8, "title": "Odd Perfect Numbers", "category_id": 1, "category": {"id": 1, "name": "number_theory", "display_name": "Number Theory", "description": "Properties of integers, prime numbers, Diophantine equations.", "slug": "number-theory", "order_index": 1, "created_at": "2026-07-31T15:26:25.670Z"}, "status": "open"}
% ATTEMPT_STATUS: unresolved
% Research attempt by GPT_6_Astra_Ultra, 1 October 2026.
% Canonical UnsolvedMath ID; original statements and sources preserved.
% FIRST_POSTED: 2026-10-01
% !TeX program = lualatex
\documentclass[11pt]{article}
\usepackage[margin=25mm]{geometry}
\usepackage{fontspec}
\setmainfont{Latin Modern Roman}
\usepackage{amsmath,amssymb,mathtools}
\usepackage{unicode-math}
\setmathfont{Latin Modern Math}
\usepackage{xurl,hyperref}
\hypersetup{colorlinks=true,urlcolor=blue}
\setcounter{secnumdepth}{0}
\setlength{\parindent}{0pt}
\setlength{\parskip}{5pt}
\setlength{\emergencystretch}{3em}
\begin{document}
\section*{\texorpdfstring{Odd Perfect Numbers}{Odd Perfect Numbers}}\label{8}
\subsection{Definitions and mathematical statement}
For an integer $N\ge1$, let $\sigma(N)=\sum_{d\mid N,\ d>0}d$. A perfect number satisfies $\sigma(N)=2N$, equivalently the sum of its positive proper divisors equals $N$. The question is
\[
 \exists N\in{\mathbb{Z}}_{>0}\quad N\equiv1\pmod2\quad\text{and}\quad\sigma(N)=2N\ ?
\]
For a prime factorization $N=\prod p_i^{a_i}$, the condition is
$\prod_i(1+p_i+\cdots+p_i^{a_i})=2\prod_i p_i^{a_i}$, with all $p_i$ odd. Either an example or an impossibility proof resolves the existence question; excluding any finite size range alone does not.
\par\smallskip\textit{Formulation and concept references:} \href{https://www.ams.org/mcom/2015-84-295/S0025-5718-2015-02941-X/S0025-5718-2015-02941-X.pdf}{[S1]}.

\subsection{Short English statement}
Can an odd positive integer equal the sum of all its positive proper divisors?
\subsection{Research attempt}

\paragraph{Scope, provenance, and selected route.}
This attempt addresses UnsolvedMath 8, with the existence quantifier exactly
as stated above. All deductions below concern a hypothetical odd perfect
integer; they neither construct one nor prove that none exists. Nielsen's
publisher abstract [E1] states a much stronger lower bound on the number
of distinct prime factors than the elementary bounds derived here. That
computation is not reproduced. The purpose here is to build an explicit
factorization obstruction, then identify why its branches do not close.

\paragraph{The parity equation forces Euler's shape.}
Write $N=\prod_{i=1}^k p_i^{a_i}$, where the $p_i$ are distinct odd primes.
Because $\sigma$ is multiplicative on coprime integers,
\[
 \prod_{i=1}^k\sigma(p_i^{a_i})=2N.
\]
For odd $p$, the sum $\sigma(p^a)=1+p+\cdots+p^a$ has parity $a+1$.
The right side has exactly one factor of two, so exactly one exponent
is odd. Denote its prime by $q$ and exponent by $\alpha$.
Moreover $\sigma(q^\alpha)$ is divisible by two but not four.
If $q\equiv3\pmod4$, grouping its even number of summands into pairs
gives terms $q^{2j}(1+q)$, each divisible by four, a contradiction.
Thus $q\equiv1\pmod4$. Now every summand is one modulo four, so
$\sigma(q^\alpha)\equiv\alpha+1\pmod4$ forces $\alpha\equiv1\pmod4$.
Consequently every candidate has the form
\[
 N=q^{4b+1}m^2,\qquad q\equiv1\pmod4,\quad (q,m)=1,
 \quad m\text{ odd},\quad b\ge0.
\]
This is a necessary condition, not a sufficient parametrization.
For example $5\cdot9=45$ has the shape but has divisor sum $78$.

\paragraph{Abundancy gives rigorous pruning inequalities.}
Set $I(t)=\sigma(t)/t$. For a prime power,
\[
 I(p^a)=1+p^{-1}+\cdots+p^{-a}
       =\frac{p}{p-1}(1-p^{-a-1})<\frac p{p-1}.
\]
Every additional positive exponent raises $I$, and coprime factors
multiply it. Hence a partially specified divisor $D\mid N$ with
$I(D)>2$ cannot extend to a perfect number. If $I(D)=2$, no proper
multiple containing an extra prime or increased exponent can be perfect.
At the other end, if a proposed full prime support $P$ satisfies
$\prod_{p\in P}p/(p-1)\le2$, every number on that support has
abundancy strictly below two. Thus no candidate has only one or two
distinct primes, since even the maximal two-prime product is
$(3/2)(5/4)=15/8<2$. If three primes all exceed three, the product is
at most $(5/4)(7/6)(11/10)=77/48<2$.

These arguments do not bound the total number of prime factors. In a
search with an imposed support-size limit, ordered remaining primes can
give useful upper bounds, but that imposed limit is an additional
hypothesis. It cannot disappear when converting a finite search into an
existence proof or a global contradiction.

\paragraph{Divisor-sum chains with their exact hypotheses.}
If $p^a\Vert N$, every odd prime divisor of $\sigma(p^a)$ must divide
$N$, because $\sigma(p^a)\mid2N$. The exact-exponent condition matters:
$p^2\mid N$ alone does not imply $\sigma(p^2)\mid\sigma(N)$.
For instance $\sigma(3^2)=13$ does not divide $\sigma(3^4)=121$.
On the branch with the indicated exponents equal to two, direct
factorization gives
\[
 \sigma(3^2)=13,\quad \sigma(13^2)=3\cdot61,\quad
 \sigma(61^2)=3\cdot13\cdot97,\quad
 \sigma(97^2)=3\cdot3169.
\]
Thus $3^2\Vert N$ forces $13\mid N$; if additionally
$13^2\Vert N$, it forces $61\mid N$, and so on. But each new prime
could have a larger even exponent, or could be the single special prime
when its residue permits. Closing the exponent-two branch does not
close these alternatives. This is the first concrete obstacle to a
naive infinite-chain argument.

There is also a valuation constraint that is stronger than simply
recording the support. For each odd prime $r\mid N$,
\[
 v_r(N)=\sum_{p^a\Vert N}v_r(\sigma(p^a)),
\]
since $v_r(2)=0$. In particular contributions to the same forced prime
must be added. A chain graph that records only whether a prime occurs
can miss contradictions caused by excessive valuation, and can equally
miss valid branches that absorb several contributions in a higher
exponent. The displayed equation is an exact conservation law; no
independence of the individual divisor-sum factors is assumed.

\paragraph{A useful two-sided constraint on the square part.}
The Euler decomposition gives
\[
 I(m^2)=\frac2{I(q^{4b+1})},\qquad
 \frac{2(q-1)}q<I(m^2)\le\frac{2q}{q+1}.
\]
The strict lower inequality uses $I(q^{4b+1})<q/(q-1)$;
the upper inequality uses its first two terms $1+1/q$.
Thus when $q$ is large the square part must have abundancy close to two
from below. This is a genuine restriction but still permits arbitrarily
many prime factors. Replacing ``close to two'' by equality, or assuming
a fixed positive distance between abundancy values and two, is invalid:
the denominators and supports are not bounded.

\paragraph{Executed checks and their coverage.}
The accompanying integer script factors every integer at most $200000$
using a smallest-prime-factor sieve and evaluates the geometric-series
formula for $\sigma$. It finds no odd perfect number in that interval
and independently checks all four displayed divisor-sum factorizations.
The files are \texttt{checks.py} and \texttt{results.json} in
\path{_Local/Erdos_problems/UnsolvedMath/attempt_audit_2026_10_01/number_theory/}.
These are small reproducibility checks, not a new record or a rerun of
Nielsen's extensive computation. The finite bound says nothing about
an odd integer above it.

\paragraph{Remaining gap and outcome.}
The attempt establishes Euler's shape, explicit abundancy exclusions,
and support and valuation propagation with the exact hypotheses visible.
A proof of nonexistence would require a termination argument covering
all possible supports and exponents, or a new global incompatibility
among these equations. Neither is obtained. In particular no estimate
here prevents branches with increasing numbers of prime divisors.
\textbf{Outcome: UNRESOLVED.} The elementary deductions are not claimed
as new results, and no finite computation is promoted to a global proof.

\subsection{Sources}
\begin{itemize}
\item[S1]Pace P. Nielsen, Mathematics of Computation 84 (2015), 2549-2567.
\url{https://www.ams.org/mcom/2015-84-295/S0025-5718-2015-02941-X/S0025-5718-2015-02941-X.pdf}.
\textit{Formulation source.}

\item[E1] Pace P. Nielsen, \emph{Odd perfect numbers, Diophantine
equations, and upper bounds}, Mathematics of Computation 84 (2015),
2549--2567. Publisher abstract inspected 1 October 2026; PDF retrieval
was unavailable. \url{https://www.ams.org/journals/mcom/2015-84-295/S0025-5718-2015-02941-X/}.
The reported ten-distinct-prime bound is background and is not a premise
of the elementary arguments above.

\end{itemize}
\end{document}
